\documentclass[
%handout,
%draft,
compress]
{beamer}

\usetheme[authorwidth=.20,titlewidth=.50,datewidth=.30,section=true,subsection=true,]{Lund}
\setbeamertemplate{navigation symbols}{}
\usefonttheme[onlymath]{serif}
% automatic table of contents at every section
\AtBeginSection[]{\frame{\tableofcontents[currentsection]}}
%\usepackage{pdfsync}
\usepackage[utf8x]{inputenc}
\SetUnicodeOption{mathletters}
\usepackage{verbatim}
\hypersetup{ % information in the PDF file.
  pdftitle={Catch and throw\\Ball and Beam},%
  pdfauthor={Sven Svensson},
%  pdfsubject={Advanced Mathematics},
%  pdfkeywords={groups, triangles, algebras}
}

\title[Short Title]{Title of that Presentation}
\author[Short Name]{Sven Svensson}

\date{2013-09-01}
\begin{document}
\frame{\titlepage}

\section*{Outline}

\part{Main Talk}
\frame{\partpage}

\frame
{
  \frametitle{Outline}
  \tableofcontents[part=1,pausesections]%,shadesubsections]
}

\section{Beamer}

\frame{\tableofcontents[current]}

\subsection{Demo}

\begin{frame}
\frametitle{Magic!}
\uncover<2->{One can display and remove objects on the same slide:}
\begin{overlayarea}{\textwidth}{6cm}
\begin{itemize}
\item<3-> \uncover<3->{{\color{structure}Words in disarray.}}
\only<4-8>{\newline \uncover<8->{That} \uncover<7->{phrase} \uncover<6->{appears} \uncover<5->{backwards} \uncover<4->{!}}
\item<9-> \uncover<9->{{\color{structure}Words that replace others.}}
\only<10-12>{\newline \only<10>{Blabla \dots}\only<11>{\dots\ blablabla \dots}\only<12>{\dots\ blablabli, \dots}}
\item<13-> \uncover<13->{{\color{structure}A Table.}}
\only<14-19>{
\begin{center}
\begin{tabular}{|c|c|c|}
\hline
\uncover<17->{AAA} & \uncover<19->{BBBB} & \uncover<15->{CCC} \\
\hline
\uncover<14->{$\mathbb D$} & \uncover<16->{$\mathbb R$} & \uncover<18->{$\mathbb F$}\\
\hline
\end{tabular}
\end{center}
}
\item<20-> \uncover<20->{{\color{structure}And many other things!}}
\begin{itemize}
\item<21->{several columns;}
\item<22->{fixed size zones;}
\item<23->{print or handout versions;}
% v \item<24->{transitions;}
\item<24->{hyperlinks;}
\item<25->{supplementary slides for questions;}
%\item<27->{works with \texttt{latex} as well as with \texttt{pdflatex}.}
\end{itemize}
\end{itemize}
\end{overlayarea}
\end{frame}

\begin{frame}
% Transition pour ce transparent, 
% la seule utilisÈe dans cette prÈsentation !
\transsplithorizontalout<1>
\frametitle{Maths with style}
Maths formula can also be dynamical!

\uncover<2->{A step by step computation\dots}
\begin{align*}
\uncover<3->{\langle d\omega, \sigma \rangle} &
\uncover<5->{\alert<5>{\ = \int_{\Delta_p} \sigma^\ast d\omega}}
\uncover<7->{\alert<7>{\ = \int_{\Delta_p} d( \sigma^\ast \omega)}} \\
 & \uncover<9->{\alert<9>{\ = \int_{\partial \Delta_p} \sigma^\ast \omega}}
 \uncover<10->{\alert<10>{\ = \sum_{i=0}^p (-1)^i \int_{\Delta_{p-1}} F^{p\ \ast}_i \sigma^\ast \omega}}\\ 
 & \uncover<10->{\alert<10>{\ =}} \uncover<8->{\ \alert<8>{\sum_{i=0}^p (-1)^i \int_{\Delta_{p-1}} (\sigma \circ F^p_i)^\ast \omega = }}
 \uncover<6->{\alert<6>{\int_{\Delta_{p-1}} (\partial \sigma )^\ast \omega}} \\
 & \uncover<4->{= \langle \omega, \partial \sigma \rangle}
\end{align*}
\end{frame}

\subsection{Fractional Derivatives}

\begin{frame}{Fractional Derivatives}
$F$ is the Fourier Transform
$$F\circ D^n = (i2\pi\xi)^n F$$
$\alt<2>{\alpha }{n}$
\begin{equation*}
D^{\alt<2>{\alpha}{n}} f
{(i2\pi k)^{\alt<2>{\alpha}{n}}}
 (i2\pi k)^{\alt<2>{\alpha}{n}} F f
\end{equation*}
\uncover<2>{$$(i2\pi k)^{\alpha} =  {|2\pi k|}^{\alpha} \exp\left(i\frac{\pi}{2}\alpha\mathrm{sgn}(k)\right)$$}
\end{frame}


\subsection{Documentation}
\begin{frame}{Documentation}
Make sure to read the excellent documentation for beamer:

\url{http://www.ctan.org/tex-archive/macros/latex/contrib/beamer/doc/beameruserguide.pdf}
\end{frame}

\begin{frame}{Columns}
You may divide the page in columns and/or blocks as follows:
\begin{columns}
\begin{column}{.4\textwidth}
\begin{block}{Left}
Left column.
\end{block}
\end{column}
\begin{column}{.4\textwidth}
\begin{block}{Right}
Right column.
\end{block}
\end{column}
\end{columns}
\end{frame}

\begin{frame}[fragile]{Fragile frames}
If you have some verbatim text in the page \emph{make sure to use the fragile option.}
\begin{verbatim}
def some_function(arg):
  print arg
\end{verbatim}
\pause
(Note that you should use syntax highlighting for the example above.)
\end{frame}

\begin{frame}[fragile]{Graphic inclusion}
To include graphics, you should simply use \verb|includegraphics| as usual:
\begin{center}
\includegraphics[height=0.7\textheight]{logoonly}
\end{center}
\end{frame}

\part{References}
\frame{\partpage}

\begin{frame}
\frametitle<presentation>{References}  
  \begin{thebibliography}{10}
  \beamertemplatearticlebibitems
  \bibitem{vonNeumann-self-adjoint}
    A.~Devinatz and A.~E.~Nussbaum and J.~von~Neumann
    \newblock On the permutability of self-adjoint operators.
    \newblock {\em Ann. of Math.} (2) {\bf 62} (1955). 199--203
    \pause
  \bibitem{vonNeumann-matrices}
    H.~H.~Goldstine and J.~von~Neumann
    \newblock Numerical inverting of matrices of high order.
    \newblock {\em Proc. Amer. Math. Soc.}, {\bf 2} (1951). 188--202
\pause
  \end{thebibliography}
\end{frame}


\end{document}

